\documentclass[10pt,a4paper]{amsart}
\usepackage[margin=19mm]{geometry}
\usepackage{amsmath,amssymb,mathtools}
\usepackage[scr=rsfs,cal=euler]{mathalfa}
\usepackage{xfrac}
\usepackage[unicode,hidelinks,bookmarks=false]{hyperref}
\newcommand{\ie}{i.e.\ }
\newcommand{\cf}{cf.\ }
\newcommand{\resp}{respectively\ }
\newcommand{\Subsection}{\subsection}
\providecommand{\nobreakdash}{\nobreak\hbox{-}}
\setlength{\emergencystretch}{2em}
\multlinegap0pt
\usepackage{bm}
\usepackage{bbm}
\usepackage{dsfont}
\usepackage{xspace}
\usepackage{cancel}
\usepackage{mathtools}
\usepackage{amsthm}
\makeatletter
\@ifundefined{definition}{\newtheorem{definition}{Definition}}{}
\@ifundefined{lemma}{\newtheorem{lemma}[definition]{Lemma}}{}
\@ifundefined{proposition}{\newtheorem{proposition}[definition]{Proposition}}{}
\makeatother
\newcommand{\ele}{\mathcal L} \newcommand{\de}{\mathcal D} \newcommand{\ka}{\mathcal K}
\providecommand{\U}[1]{\protect\rule{.1in}{.1in}}
\title[Notes on Leibniz $n$-algebras]{Notes on Leibniz $n$-algebras}
\author{Jos\'e Manuel Casas, Emzar Khmaladze, Manuel Ladra}
\date{Source: arXiv:2601.17531v1 (CC BY 4.0)}
\begin{document}
\maketitle

\noindent\emph{Source and licence.} Notes on Leibniz $n$-algebras. Version arXiv:2601.17531v1 (CC BY 4.0), distributed under the Creative Commons Attribution 4.0 International licence, as recorded in the arXiv licence metadata for this exact version and shown on the abstract page https://arxiv.org/abs/2601.17531v1. Full rights notice: licenses/arxiv-2601.17531-CC-BY-4.0.txt.\par\smallskip

\noindent\textit{Editorial scope.} Counted unit: the Lemma whose source label is \texttt{Lemma action} in the article (source0.tex lines 364--366, source label \texttt{Lemma action}), delivered together with the article's own complete original proof of it (source0.tex lines 367--382, one \texttt{proof} environment of 16 lines). No mathematics is written, repaired or re-derived: statement and proof are the article's own text line for line. Required context retained uncounted: FI (lines 153--162); leib (lines 283--291); Def action (lines 342--359). Every definition, display and label of those lines is retained unchanged. Omitted from this package and not redistributed: 1--152, 163--282, 292--341, 360--363, 383--925. Nothing inside a retained excerpt is omitted or altered: every retained body is delivered exactly as the article's lines are written, so no omission receipt of any kind accompanies this package. Citations: the retained excerpts carry no citation command at all, so no bibliography entry of the article is transcribed and no reference is redistributed. Reference adaptation: none was needed and none was made; every reference inside the delivered unit resolves inside it, so no unresolved reference or citation is produced. Printed slips kept literal: no printed word, symbol, hypothesis or number of the counted statement or of its proof was altered. Layout and class: the article's own class is \texttt{amsart} with packages including amssymb, amsmath, amsfonts, fontenc, inputenc, xy, xcolor, amsthm, a4wide, mathtools, orcidlink, multicol, hyperref, graphicx; the wrapper is the lane's pinned \texttt{amsart} editorial wrapper at 10pt on A4, which restates the theorem-like environments the retained text needs and adds the article's own macro definitions that the retained text needs (\texttt{\textbackslash U}, \texttt{\textbackslash ele}), with extra wrapper packages (bm, bbm, dsfont, xspace, cancel, mathtools). The delivered numbering is the unit's own. Assets: no retained span contains a figure environment, a graphics-inclusion command, a PDF-page inclusion, a table or a TikZ picture; no image file is copied into this package, the package declares no asset dependency and no image file is redistributed with it. Licence: version arXiv:2601.17531v1 of this article is distributed under the Creative Commons Attribution 4.0 International licence, as recorded in the arXiv licence metadata for this exact version and shown on the abstract page https://arxiv.org/abs/2601.17531v1. A full rights notice is delivered at licenses/arxiv-2601.17531-CC-BY-4.0.txt. Characters: every retained excerpt is pure ASCII, so the delivered engine, XeLaTeX with its default Latin Modern text font, reports no missing character.
\begin{definition}  A Leibniz $n$-algebra  is a vector space ${\ele}$
equipped with an $n$-linear operation ($n$-ary bracket) $[-, \dots,-] : {\ele}^{\otimes n} \to {\ele}$  satisfying the following fundamental identity for all $x_1, \dots, x_n, y_2, \dots, y_n \in {\ele}$:
\begin{equation} \label{FI}
\begin{array}{l}
[[x_1, x_2, \dots, x_n],y_2,\dots,y_{n}] =
  \displaystyle \sum_{i=1}^n [x_1,\dots, x_{i-1}, [ x_i,y_2,\dots,
y_{n}],x_{i+1},\dots,x_n]
\end{array}
\end{equation}
\end{definition}

\begin{proposition} \label{leib}
Let ${\ele}$ be a Leibinz $n$-algebra with respect to the $n$-ary bracket \linebreak $[-,\stackrel{n}\dots,-] : {\ele}{^{\otimes n}} \to \ele$. Then ${\ele}$ is a Leibniz $p$-algebra  with respect to the $p$-ary bracket
\[
[l_1,l_2,\dots,l_{p}]
  =[l_1, \dots l_{n-1},[l_n, \dots, l_{2n-2}, [\dots,[l_{p-n+1}, \dots, l_{p}] \dots]]].
\]
 %$\omega_p :
%{\ele}^{\otimes p} \to {\ele}$ given by (\ref{structure}).
\end{proposition}

\begin{definition}\label{Def action} We say that a Leibniz $n$-algebra $\mathcal{P}$ acts on
a Leibniz $n$-algebra $\mathcal{L}$ if $2^n-2$ $n$-linear maps
\[
[-,\stackrel{n}{\cdots},-]:\mathcal{L}^{\otimes i_1}\otimes \mathcal{P}^{\otimes
j_1}\otimes\cdots\otimes \mathcal{L}^{\otimes i_m}\otimes
\mathcal{P}^{\otimes j_m}\to \mathcal{L}
\]
are given, where $1\leq m\leq n-1$, $\underset{\scriptsize{1\leq
k\leq m}}\sum (i_k+j_k)=n$, $0\leq i_k \leq n-1$ and at least one
$i_k\neq 0$, $0\leq j_k \leq n-1$ and at least one $j_k\neq 0$,
such that $2^{2n-1}-2$ equalities hold which are obtained from the fundamental identity
(\ref{FI}) by taking exactly $i$ of the variables
$x_1,\dots,x_n,y_1,\dots, y_{n-1}$ in $\mathcal{L}$ and all the
others in $\mathcal{P}$, %($\binom{2n-1}{i}$ equalities)
and by changing $i=1,\dots,2n-2$.

%The action is called trivial if all these linear maps are trivial.
\end{definition}

\begin{lemma}\label{Lemma action}
Let $\mathcal{P}$ and $\mathcal{L}$ be Leibniz $n$-algebras and $\mathcal{P}$ acts on $\mathcal{L}$. Then the  Leibniz $p$-algebra $\U_n^p(\mathcal{P})$ acts on the Leibniz $p$-algebra $\U_n^p(\mathcal{L})$.
\end{lemma}

\begin{proof}
Since $\U_n^p(\mathcal{P})=\mathcal{P}$ and $\U_n^p(\mathcal{L})= \mathcal{L} $ are Leibniz $p$-algebras with the $p$-ary brackets as described in Proposition \ref{leib}, we need to construct $2^p-2$ $p$-linear maps
\[
[-,\stackrel{p}{\cdots},-]:\mathcal{L}^{\otimes i_1}\otimes \mathcal{P}^{\otimes
j_1}\otimes\cdots\otimes \mathcal{L}^{\otimes i_m}\otimes
\mathcal{P}^{\otimes j_m}\to \mathcal{L},
\]
with $\underset{\scriptsize{1\leq
k\leq m}}\sum (i_k+j_k)=p$ and satisfying the same conditions as in Definition \ref{Def action}. Let us fix such $i_1, \cdots, i_m$, $j_1, \cdots, j_m$
and define the corresponding $p$-linear map $[-,\stackrel{p}{\cdots},-]$ by
\begin{align}\label{equation action}
[x_1, & \dots, x_{i_1}, x_{i_1 +1}, \dots, x_{i_1+j_1}, \dots,   x_{p}] \notag \\
 & =[x_1, \dots x_{n-1},[x_n, \dots, x_{2n-2}, [\dots,[x_{p-n+1}, \dots, x_{p}] \dots]]],
\end{align}
where any $x$ element belongs to either $\mathcal{L}$ or $\mathcal{P}$, but at the same time, at least one of these elements belongs to $\mathcal{L}$ and at least one belongs to $\mathcal{P}$. In the right hand side of equation (\ref{equation action}) the iterated $n$-ary brackets are given by the action of $\mathcal{P}$ on $\mathcal{L}$ and therefore satisfy the fundamental equation (\ref{FI}). Then, by the same argument as in Proposition \ref{leib}, it is easy to see that $p$-linear brackets defined by (\ref{equation action}) also satisfy the corresponding fundamental identity, i.e. they indeed define an action of $\U_n^p(\mathcal{P})$ on $\U_n^p(\mathcal{L})$.
\end{proof}

\end{document}
